\section{Introduction}
\label{sec:introduction}

The abundance conjecture predicts that a nef log canonical divisor
$K_X+B$ on a projective log canonical pair is semiample: some positive
Cartier multiple is generated by its global sections. Semiampleness
turns the numerical positivity of the divisor into a morphism. We
consider threefolds in positive characteristic and the case of numerical
dimension one. For a nef $\Q$-Cartier divisor $L$ on a threefold, this
condition means that, for an ample Cartier divisor $H$,
\[
LH^2>0\qquad\text{and}\qquad L^2H=0.
\]

\begin{theorem}\label{thm:main}
Let $k$ be an algebraically closed field of characteristic $p>3$.
Let $(X,B)$ be a projective log canonical threefold pair over $k$,
where $X$ is normal, $B$ is an effective $\Q$-divisor, and $K_X+B$
is $\Q$-Cartier. If $K_X+B$ is nef and has numerical dimension one,
then $K_X+B$ is semiample.
\end{theorem}

The boundary coefficients are arbitrary rational numbers allowed by
log canonicity, and neither terminality nor $\Q$-factoriality is assumed
for the original variety. The theorem proves the numerical-dimension-one
case of threefold log abundance in this characteristic range.

Over the complex numbers, Miyaoka proved abundance for minimal
threefolds of numerical dimension one \cite{Miyaoka}; Kawamata completed
abundance for minimal threefolds and gave another proof of this case
\cite[Section~4]{Kawamata}. Keel, Matsuki, and McKernan established
log abundance for log canonical threefold pairs \cite{KMM,KMMCorrection}. In positive
characteristic, Keel's semiampleness criterion \cite{Keel} and the
basepoint-freeness methods of Cascini, Tanaka, and Xu \cite{CTX}
provided foundations. The threefold minimal model program
developed through the work of Hacon--Xu, Birkar, and Birkar--Waldron in
characteristic greater than five \cite{HX,Birkar,BW}, the log canonical
and relative extensions of Waldron and Hashizume--Nakamura--Tanaka
\cite{Waldron,HNT}, and Hacon--Witaszek's treatment of characteristic
five \cite{HW}.

For abundance, Xu--Zhang proved nonvanishing for terminal minimal
threefolds in characteristic greater than five \cite{XZ}. Witaszek
obtained klt nonvanishing and abundance in nef dimension at most two
in that characteristic range \cite{Witaszek}. The nef dimension $n(L)$
is the dimension of the nef reduction of $L$. On a threefold over an
uncountable algebraically closed field, $n(L)=3$ means that no curve of $L$-degree zero passes
through a very general point. Zheng Xu established nonvanishing for pseudoeffective log
canonical divisors on projective lc threefold pairs in characteristic
greater than three, together with
abundance in nef dimension at most two
\cite[Theorems~1.4 and~1.8]{XuNonvanishing}; he subsequently proved
abundance in numerical dimension two \cite[Theorem~1.1]{XuNuTwo}.
For the present problem, the remaining case has numerical dimension
one and nef dimension three.

The companion paper \cite{Companion} treats terminal threefolds with
nef canonical divisor of numerical dimension one. Its proof isolates
a stronger structural result: certain connected nef Cartier divisors
on smooth threefolds cannot exist. We state this cycle obstruction
in Theorem~\ref{thm:obstruction}, with all its hypotheses. The present
paper constructs such a divisor from a log canonical pair with
arbitrary rational boundary. The companion's obstruction is an
essential input to the proof.

\subsection*{The boundary and the proof strategy}

After an uncountable field extension and a crepant birational change,
we may suppose that the underlying variety is terminal and
$\Q$-factorial. Nonvanishing gives a nonzero effective Cartier divisor
$N_0\sim m(K_X+B)$. We assume for contradiction that $n(N_0)=3$.
The main additional issue is to obtain a relation for a power of the
canonical bundle near a component of this \emph{log} canonical divisor.
The boundary and the exceptional divisors of a resolution must both
be controlled there.

Proposition~\ref{prop:separation} supplies that control for the boundary:
it proves $BN_0H=0$ and shows that every boundary component is either
contained in $\Supp N_0$ or disjoint from it. If the intersection were
positive, complete-intersection curves would have arbitrarily small
ratios of $N_0$-degree to anticanonical degree. Quantitative
bend-and-break \cite[Theorem~1.1]{JLR} would then produce
$N_0$-trivial rational curves through very general points. This
separation argument is formulated independently of the MMP.

Next we adapt the isolation construction of
\cite[Section~2.3]{Companion}. Raising the boundary to its log canonical
threshold along $N_0$ gives a coefficient-one component on a dlt model.
A partial MMP, trivial for the resulting full log canonical class,
makes this component disjoint from the rest of the divisor.
We use Xu's partial-MMP formulation
\cite[Section~2.3]{XuNuTwo}, whose ray argument adapts
\cite[Lemma~5.1]{KMM}. Reduced-boundary adjunction and surface
abundance \cite{Tanaka,Posva} then yield torsion on a Cartier
multiple of that component. We include the companion's extension
argument across finitely many points, since the required torsion is
on the full nonreduced Cartier scheme.

Finally, terminality makes the singular locus finite, so the exceptional
divisors of a common resolution obey the same support separation as
the boundary. A finite separable cover removes the prime-to-$p$ part
of the pluricanonical exponent. A connected component of the pulled-back
Cartier divisor, divided by the gcd of its multiplicities, has exactly
the numerical, torsion, and canonical properties forbidden by
Theorem~\ref{thm:obstruction}.

Section~\ref{sec:foundations} states the obstruction and makes the
initial reductions. Section~\ref{sec:boundary} proves boundary
separation, and Section~\ref{sec:isolation} isolates a component.
Section~\ref{sec:torsion} proves torsion on its Cartier thickening.
Section~\ref{sec:cycle} constructs the final cycle and completes the
proof.
